% Generated draft: requires mathematical and source review.
\documentclass[11pt]{article}
\usepackage[margin=25mm]{geometry}
\usepackage{fontspec}
\setmainfont{Latin Modern Roman}
\usepackage{amsmath,amssymb,mathtools,unicode-math}
\setmathfont{Latin Modern Math}
\usepackage{enumitem,needspace,xurl,hyperref}
\setlength{\parindent}{0pt}
\setlength{\parskip}{5pt}
\setlength{\emergencystretch}{3em}
\newcommand{\N}{\mathbb{N}}
\newcommand{\Z}{\mathbb{Z}}
\newcommand{\Q}{\mathbb{Q}}
\newcommand{\R}{\mathbb{R}}
\newcommand{\C}{\mathbb{C}}
\newcommand{\F}{\mathbb{F}}
\newcommand{\A}{\mathbb{A}}
\newcommand{\PP}{\mathbb{P}}
\newcommand{\E}{\mathbb{E}}
\newcommand{\eps}{\varepsilon}
\newcommand{\dd}{\,\mathrm{d}}
\begin{document}
\section*{UnsolvedMath ID: 30002142}
\textbf{Source problem code:} catalogue entry
\subsection{Definitions and mathematical statement}
A contact form on a smooth manifold \(M\) of dimension \(2n+1\) is a smooth one-form \(\alpha\) satisfying \(\alpha\wedge(\dd\alpha)^n\neq 0\) everywhere. Write \(\xi=\ker\alpha\) for its contact distribution. Define the strict contactomorphism group and the contactomorphism group by
\[
\operatorname{Cont}(M,\alpha)=\{f\in\operatorname{Diff}(M):f^*\alpha=\alpha\},
\qquad
\operatorname{Cont}(M,\xi)=\{f\in\operatorname{Diff}(M):f_*\xi=\xi\}.
\]
Once topologies on these groups have been specified, the subscript \(0\) denotes the connected component of the identity in the respective group. Since preserving \(\alpha\) implies preserving \(\xi\), the inclusion induces
\[
j:\operatorname{Cont}_0(M,\alpha)\hookrightarrow\operatorname{Cont}_0(M,\xi),\qquad j(f)=f.
\]
\textbf{Source question.} For a contact manifold equipped with a contact form \(\alpha\), determine the topology of this natural inclusion \(j\).

\textbf{Unresolved qualifications.} This is an open-ended request to describe an inclusion, rather than a specified proposition about a particular invariant. The supplied statement does not specify the topology on the transformation groups; a customary interpretation uses a smooth \(C^\infty\) topology, but its precise meaning must be fixed, particularly for noncompact \(M\). It also does not specify dimension restrictions, compactness, boundary conditions, support conditions, or whether \(M\) is assumed connected. The notation \(\alpha\) requires a globally specified contact form. The definition of \(\operatorname{Cont}(M,\xi)\) above allows all distribution-preserving diffeomorphisms; the source does not explicitly state a coorientation convention. Finally, although the supplied literature assessment discusses general regular contact manifolds, the displayed question does not impose regularity. The intended scope and the precise topological information requested therefore require review.
\subsection{Short English statement}
The question asks how the identity component of the group preserving a contact form sits inside the identity component of the group preserving its contact distribution. The group topology, manifold hypotheses, and intended topological invariants are not specified. The supplied dated curation labels the problem partially solved, describing progress for the standard three-sphere and an unresolved general regular case; these are supplied assessments requiring independent review.
\subsection{Sources}
\begin{itemize}
\item Local UnsolvedMath registry, permanent ID \(30002142\), source problem code \texttt{catalogue entry}, title ``Topology of Strict-to-General Contactomorphism Inclusions.'' The record gives proposed year 2012 and names no proposer.
\item Supplied report citation: \emph{Geometric Group Theory, Hyperbolic Dynamics and Symplectic Geometry} (2013), Oberwolfach Reports, pp.\ 2139--2203, DOI \texttt{10.4171/owr/2012/35}. Also listed as ``Oberwolfach Report 35'': \url{https://doi.org/10.4171/owr/2012/35}. The background attributes the question to workshop open-problem discussions documented in Oberwolfach Reports.
\item Supplied literature reference: \emph{Chern--Weil theory and strict contactomorphisms}, \url{https://arxiv.org/abs/1409.1913}.
\item Literature assessment dated exactly \texttt{2026-08-21}: the supplied curation labels the problem ``partially solved'' and states that the standard three-sphere case has a strong homotopy-theoretic description of the strict-contact subgroup inclusion, while the topology of the inclusion for general regular contact manifolds remains unresolved. This is supplied dated curation derived from the OWR material, not an independently verified assessment; time-sensitive status and cited claims require independent verification.
\item The registry labels the statement \texttt{corrected\_verified} and reports that its mathematical content was checked against the official Oberwolfach report, with notation, context, and question form normalized without changing the source-level claim. This verification is an assertion of the supplied record and was not independently repeated here.
\end{itemize}
\end{document}
